\documentclass[12pt,a4paper]{article}

% ── Packages ─────────────────────────────────────────────────────────────────
\usepackage[margin=2.5cm]{geometry}
\usepackage{amsmath,amssymb}
\usepackage{graphicx}
\usepackage{booktabs}
\usepackage{array}
\usepackage{xcolor}
\usepackage{hyperref}
\usepackage{float}
\usepackage{caption}
\usepackage{subcaption}
\usepackage{enumitem}
\usepackage{parskip}
\usepackage{titlesec}
\usepackage{fancyhdr}
\usepackage{listings}
\usepackage{microtype}

% ── Colours ───────────────────────────────────────────────────────────────────
\definecolor{alexblue}{RGB}{0,62,120}
\definecolor{lightgray}{RGB}{240,240,240}
\definecolor{codegray}{RGB}{245,245,245}

% ── Hyperref setup ────────────────────────────────────────────────────────────
\hypersetup{
  colorlinks=true,
  linkcolor=alexblue,
  urlcolor=alexblue,
  citecolor=alexblue,
}

% ── Section formatting ────────────────────────────────────────────────────────
\titleformat{\section}
  {\large\bfseries\color{alexblue}}{\thesection}{1em}{}[\titlerule]
\titleformat{\subsection}
  {\normalsize\bfseries\color{alexblue}}{\thesubsection}{1em}{}

% ── Header / Footer ───────────────────────────────────────────────────────────
\pagestyle{fancy}
\fancyhf{}
\rhead{\small CSE: Pattern Recognition -- Assignment \#1}
\lhead{\small Alexandria University}
\cfoot{\thepage}
\renewcommand{\headrulewidth}{0.4pt}

% ── Code listing style ────────────────────────────────────────────────────────
\lstset{
  backgroundcolor=\color{codegray},
  basicstyle=\ttfamily\footnotesize,
  breaklines=true,
  frame=single,
  rulecolor=\color{lightgray},
}

% ─────────────────────────────────────────────────────────────────────────────
\begin{document}

% ── Title page ────────────────────────────────────────────────────────────────
\begin{titlepage}
  \centering
  \vspace*{2cm}
  {\Large\bfseries\color{alexblue} Alexandria University\\[0.3em]
   Faculty of Engineering\\[0.3em]
   Computer and Systems Engineering Department\par}
  \vspace{1.5cm}
  \rule{\linewidth}{1.5pt}\\[0.4em]
  {\LARGE\bfseries Assignment \#1\\[0.4em]
   Multi-class Diabetes Risk Prediction\par}
  \rule{\linewidth}{1.5pt}
  \vspace{1.5cm}
  {\large\bfseries Module: K-Nearest Neighbours Classifier\par}
  \vspace{2cm}
  \begin{tabular}{ll}
    \textbf{Course:}    & CSE -- Pattern Recognition \\[0.4em]
    \textbf{Assigned:}  & Wednesday, 4 March 2026 \\[0.4em]
    \textbf{Due:}       & Wednesday, 18 March 2026 \\[0.4em]
    \textbf{Instructor:}& Prof.\ Dr.\ Marwan Torki \\[0.4em]
    \textbf{TA:}        & Eng.\ Nour Eddine \\
  \end{tabular}
  \vfill
  {\small \today}
\end{titlepage}

\tableofcontents
\newpage

% ─────────────────────────────────────────────────────────────────────────────
\section{Introduction}
% ─────────────────────────────────────────────────────────────────────────────

This report presents the K-Nearest Neighbours (KNN) component of Assignment~\#1, which
addresses multi-class diabetes risk prediction using the
\textit{Diabetes Health Indicators Dataset} (BRFSS~2015).  
The dataset contains 21~patient attributes (e.g.\ BMI, age, high blood pressure,
physical activity) and a three-class target variable:
\textbf{0}~(No Diabetes), \textbf{1}~(Prediabetes), and \textbf{2}~(Diabetes).
A central challenge is severe class imbalance: the Prediabetes class is
substantially smaller than the other two.

The goals of this module are to:
\begin{enumerate}[label=(\roman*)]
  \item tune a KNN classifier on the validation set across multiple values of $k$ and
        two distance metrics;
  \item train a baseline model on the original imbalanced data and a second model after
        applying Synthetic Minority Oversampling (SMOTE);
  \item evaluate both models on the held-out test set using accuracy, micro-, macro-,
        and weighted F1-scores; and
  \item analyse the confusion matrices and the impact of SMOTE on minority-class recall.
\end{enumerate}

% ─────────────────────────────────────────────────────────────────────────────
\section{Data Preparation}
% ─────────────────────────────────────────────────────────────────────────────

The full data pipeline is implemented in \texttt{data.py}. The key steps are
summarised below.

\subsection{Train / Validation / Test Split}

The dataset is partitioned with stratified sampling to preserve the original
class distribution in each subset:

\begin{center}
\begin{tabular}{lrr}
  \toprule
  \textbf{Split} & \textbf{Fraction} & \textbf{Approx.\ samples} \\
  \midrule
  Training   & 70\% & 177{,}580 \\
  Validation & 10\% & 25{,}369  \\
  Test       & 20\% & 50{,}736  \\
  \bottomrule
\end{tabular}
\end{center}

A fixed random seed of \texttt{42} is used throughout to guarantee
reproducibility, as required by the assignment.

\subsection{Feature Preprocessing}

All features are standardised with \texttt{StandardScaler} (zero mean,
unit variance) fitted exclusively on the training partition.  
This step is critical for KNN because the algorithm is distance-based; without
scaling, high-variance features such as BMI would dominate the distance
computation.

\subsection{Feature Selection}

To reduce dimensionality and remove uninformative features, \texttt{SelectKBest}
with \texttt{mutual\_info\_classif} is used to retain the top $K = 15$ features
out of the original~21.  
The selector is fitted on the training set only to prevent data leakage, and the
learned mask is then applied identically to the validation and test sets.

\subsection{Class Imbalance}

The original class distribution is heavily skewed towards No Diabetes.  
The Prediabetes class (class~1) accounts for less than~2\% of all samples, making
this a severe imbalance problem.  
For the \textit{balanced} variant of KNN, Synthetic Minority Oversampling
Technique (SMOTE) is applied to the training set after feature selection, generating
synthetic interpolated samples for the minority classes until all three classes are
equally represented.

% ─────────────────────────────────────────────────────────────────────────────
\section{KNN Implementation and Hyperparameter Tuning}
% ─────────────────────────────────────────────────────────────────────────────

\subsection{Algorithm Overview}

The K-Nearest Neighbours classifier is a non-parametric, instance-based learning
algorithm.  Given a query point $\mathbf{x}$, it identifies the $k$ training
samples closest in feature space and assigns the majority class label:

\[
  \hat{y}(\mathbf{x}) = \underset{c}{\arg\max}
  \sum_{i \in \mathcal{N}_k(\mathbf{x})} \mathbf{1}[y_i = c],
\]

where $\mathcal{N}_k(\mathbf{x})$ is the set of $k$ nearest neighbours under a
chosen distance metric.

Two distance metrics were evaluated:

\begin{itemize}
  \item \textbf{Euclidean} ($L_2$):
    $d(\mathbf{x},\mathbf{z}) = \sqrt{\sum_j (x_j - z_j)^2}$
  \item \textbf{Manhattan} ($L_1$):
    $d(\mathbf{x},\mathbf{z}) = \sum_j |x_j - z_j|$
\end{itemize}

\subsection{Hyperparameter Search}

The search grid spans $k \in \{3, 5, 11, 21\}$ and both distance metrics.
For each combination, a KNN model is trained on the appropriate training set
(imbalanced for the baseline, SMOTE-resampled for the balanced variant) and
evaluated on the validation set using \textbf{macro F1-score} -- the most
informative metric for an imbalanced dataset (see Section~\ref{sec:metric}).

\paragraph{Baseline search (imbalanced data).}

\begin{center}
\begin{tabular}{cccr}
  \toprule
  \textbf{Method} & $k$ & \textbf{Metric} & \textbf{Val.\ Macro F1} \\
  \midrule
  none & 3  & euclidean & 0.3935 \\
  none & 3  & manhattan & \textbf{0.3955} \\
  none & 5  & euclidean & 0.3904 \\
  none & 5  & manhattan & 0.3917 \\
  none & 11 & euclidean & 0.3784 \\
  none & 11 & manhattan & 0.3816 \\
  none & 21 & euclidean & 0.3763 \\
  none & 21 & manhattan & 0.3728 \\
  \bottomrule
\end{tabular}
\captionof{table}{Baseline (imbalanced) KNN hyperparameter grid on validation set.
  Best configuration highlighted in bold.}
\end{center}

\paragraph{SMOTE search (balanced data).}

\begin{center}
\begin{tabular}{cccr}
  \toprule
  \textbf{Method} & $k$ & \textbf{Metric} & \textbf{Val.\ Macro F1} \\
  \midrule
  smote & 3  & euclidean & 0.3998 \\
  smote & 3  & manhattan & 0.4036 \\
  smote & 5  & euclidean & 0.3971 \\
  smote & 5  & manhattan & \textbf{0.4065} \\
  smote & 11 & euclidean & 0.3894 \\
  smote & 11 & manhattan & 0.4041 \\
  smote & 21 & euclidean & 0.3837 \\
  smote & 21 & manhattan & 0.4007 \\
  \bottomrule
\end{tabular}
\captionof{table}{SMOTE-balanced KNN hyperparameter grid on validation set.
  Best configuration highlighted in bold.}
\end{center}

\paragraph{Selected configurations.}
\begin{itemize}
  \item \textbf{Baseline:} $k = 3$, Manhattan distance
        (val.\ macro F1 $= 0.3955$).
  \item \textbf{Balanced (SMOTE):} $k = 5$, Manhattan distance
        (val.\ macro F1 $= 0.4065$).
\end{itemize}

Both searches consistently favour Manhattan distance over Euclidean and smaller
values of $k$, suggesting that the decision boundaries in this 15-dimensional
scaled feature space benefit from finer local structure.  Larger $k$ values lead
to over-smoothing, blurring the boundaries around the minority Prediabetes class.

% ─────────────────────────────────────────────────────────────────────────────
\section{Evaluation}
\label{sec:eval}
% ─────────────────────────────────────────────────────────────────────────────

\subsection{Metrics Reported}

Four metrics are computed on the held-out test set:

\begin{itemize}
  \item \textbf{Accuracy:} fraction of correctly classified samples.
  \item \textbf{Micro F1:} globally aggregated precision and recall across all
        instances; equivalent to accuracy in a single-label setting.
  \item \textbf{Macro F1:} unweighted mean of per-class F1-scores; each class
        contributes equally regardless of size.
  \item \textbf{Weighted F1:} mean of per-class F1-scores weighted by class support;
        penalises poor performance on large classes more than on small ones.
\end{itemize}

\subsection{Which Metric Is Most Meaningful?}
\label{sec:metric}

For this severely imbalanced three-class problem, \textbf{macro F1-score} is the
most informative metric.  Accuracy and micro F1 are dominated by the majority class
(No Diabetes, $\approx 84\%$ of samples) and can appear deceptively high even when
the model entirely ignores Prediabetes.  Weighted F1 down-weights the minority
class contributions proportionally to support, again concealing poor minority-class
performance.  Macro F1, by treating every class equally, directly penalises the
model when it fails to recognise Prediabetes or Diabetes cases -- the clinically
critical misclassifications.

\subsection{Test Set Results}

\begin{center}
\begin{tabular}{lcccc}
  \toprule
  \textbf{Model} & \textbf{Accuracy} & \textbf{Micro F1}
                 & \textbf{Macro F1} & \textbf{Weighted F1} \\
  \midrule
  KNN Baseline (Imbalanced) & 0.8215 & 0.8215 & 0.3992 & 0.7989 \\
  KNN Balanced (SMOTE)      & 0.7080 & 0.7080 & 0.4101 & 0.7459 \\
  \bottomrule
\end{tabular}
\captionof{table}{KNN test-set performance summary.}
\end{center}

\paragraph{Per-class breakdown.}

\begin{center}
\begin{tabular}{lcccc}
  \toprule
  \textbf{Class} & \textbf{Precision} & \textbf{Recall}
                 & \textbf{F1} & \textbf{Support} \\
  \midrule
  \multicolumn{5}{l}{\textit{KNN Baseline (Imbalanced)}} \\
  No Diabetes  & 0.87 & 0.94 & 0.90 & 42{,}741 \\
  Prediabetes  & 0.05 & 0.00 & 0.01 &    926 \\
  Diabetes     & 0.37 & 0.24 & 0.29 &  7{,}069 \\
  \midrule
  \multicolumn{5}{l}{\textit{KNN Balanced (SMOTE)}} \\
  No Diabetes  & 0.91 & 0.76 & 0.82 & 42{,}741 \\
  Prediabetes  & 0.03 & 0.08 & 0.04 &    926 \\
  Diabetes     & 0.29 & 0.50 & 0.37 &  7{,}069 \\
  \bottomrule
\end{tabular}
\captionof{table}{Per-class classification report for both KNN variants.}
\end{center}

% ─────────────────────────────────────────────────────────────────────────────
\section{Confusion Matrices}
% ─────────────────────────────────────────────────────────────────────────────

Figure~\ref{fig:cm} shows the row-normalised confusion matrices for both KNN
variants on the test set.  Raw sample counts are overlaid in grey.

\begin{figure}[H]
  \centering
  \includegraphics[width=\linewidth]{confusion_matrices_knn.png}
  \caption{Row-normalised confusion matrices for KNN Baseline (left) and
           KNN Balanced -- SMOTE (right) on the test set.
           Cell values are row-normalised probabilities; grey annotations give
           raw counts.}
  \label{fig:cm}
\end{figure}

% ─────────────────────────────────────────────────────────────────────────────
\section{Analysis and Discussion}
% ─────────────────────────────────────────────────────────────────────────────

\subsection{Impact of SMOTE on Performance}

Comparing the two variants reveals a clear trade-off between overall accuracy and
minority-class recall:

\begin{itemize}
  \item \textbf{Accuracy / Micro F1} drops from 0.8215 to 0.7080 after SMOTE.
        This is expected: the model sacrifices some correct predictions for No
        Diabetes in order to improve recall for the minority classes.

  \item \textbf{Macro F1} \textit{improves} slightly from 0.3992 to 0.4101.
        Although the gain is modest, it reflects a real improvement in the model's
        ability to recognise Prediabetes and Diabetes cases.

  \item \textbf{Diabetes recall} improves substantially from 0.24 to 0.50, meaning
        the SMOTE model correctly identifies twice as many diabetic patients.
        This is clinically significant.

  \item \textbf{Prediabetes recall} improves from a near-zero~0.00 to 0.08.
        The absolute level is still very low due to the extreme rarity of this class
        ($<2\%$ of samples), but the direction of change is correct.

  \item \textbf{No Diabetes recall} decreases from 0.94 to 0.76.  
        A fraction of the majority class is now misclassified as Diabetes or
        Prediabetes -- the cost of improved minority sensitivity.
\end{itemize}

\subsection{Misclassification Patterns}

From Figure~\ref{fig:cm}, several patterns are evident:

\paragraph{Baseline model.}
The Prediabetes row is almost entirely classified as No Diabetes (0.84) or Diabetes
(0.15), with only 3 out of 926 Prediabetes instances correctly identified.  
The model has effectively learnt to ignore Prediabetes due to its negligible share
of the training data.  Diabetes recall is also modest at 0.24, with 76\% of diabetic
patients misclassified as No Diabetes.

\paragraph{Balanced (SMOTE) model.}
Prediabetes confusion is redistributed: 53\% are misclassified as No Diabetes and
39\% as Diabetes.  The diagonal entry rises from 0.00 to 0.08, indicating the model
has learned at least some Prediabetes signal.  Diabetes recall jumps to 0.50, but
a large proportion (0.41) of diabetic patients are still predicted as No Diabetes.

\paragraph{Clinical interpretation.}
The systematic confusion between Prediabetes and No Diabetes arises partly because
prediabetes is a \emph{transitional} state; its physiological markers (e.g.\ slightly
elevated BMI or borderline blood pressure) overlap substantially with the healthy
class.  Mathematically, KNN's local neighbourhood in this 15-dimensional feature
space tends to be dominated by No Diabetes neighbours due to its overwhelming
prevalence, even after SMOTE resampling of the training set.

\subsection{Strengths and Limitations of KNN}

\begin{itemize}
  \item \textbf{Strengths:} no parametric assumptions; directly captures non-linear
        decision boundaries; easy to interpret ($k$ nearest examples).
  \item \textbf{Limitations:} inference time scales with training set size ($O(n)$
        per query), making it expensive on the full BRFSS dataset;
        susceptible to the curse of dimensionality; does not naturally model
        posterior class probabilities; highly sensitive to irrelevant features even
        after feature selection.
\end{itemize}

% ─────────────────────────────────────────────────────────────────────────────
\section{Conclusion}
% ─────────────────────────────────────────────────────────────────────────────

The KNN classifier was tuned over a grid of $k \in \{3, 5, 11, 21\}$ and two
distance metrics using macro F1-score on the validation set.  The best
configuration for the baseline was $k=3$, Manhattan distance, and for the
SMOTE-balanced variant $k=5$, Manhattan distance.

SMOTE modestly improved macro F1 from 0.3992 to 0.4101 and substantially improved
Diabetes recall (0.24 $\to$ 0.50), demonstrating its value for surfacing
minority-class patterns.  However, Prediabetes recall remains very low in both
variants due to the extreme class imbalance and the high degree of feature-space
overlap with the No Diabetes class.

These results motivate the subsequent evaluation of more expressive models
(Softmax Regression and Feedforward Neural Network) which may better separate the
overlapping class regions through learned non-linear transformations.

\end{document}
